\chapter*{出版信息}
\markboth{出版信息}{出版信息}
\addcontentsline{toc}{chapter}{出版信息}

Herbert Amann

Joachim Escher

Birkhäuser Verlag Basel · Boston · Berlin

作者：

Herbert Amann

Institut für Mathematik

Universität Zürich

Winterthurerstr. 190

CH-8057 Zürich

电子邮箱：amann@math.unizh.ch

Joachim Escher

Institut für Angewandte Mathematik

Universität Hannover

Welfengarten 1

D-30167 Hannover

电子邮箱：escher@ifam.uni-hannover.de

德文原版以同一书名由瑞士 Birkhäuser Verlag 出版，© 1998 Birkhäuser Verlag

2000 数学学科分类 26-01, 26Axx; 03-01, 30-01, 40-01, 54-01

本书的 CIP 编目记录可从美国华盛顿特区国会图书馆获得。

Bibliografische Information Der Deutschen Bibliothek

Die Deutsche Bibliothek verzeichnet diese Publikation in der Deutschen Nationalbibliografie; detaillierte bibliografische Daten sind im Internet über <http://dnb.ddb.de> abrufbar.

ISBN 3-7643-7153-6 Birkhäuser Verlag, Basel – Boston – Berlin

本作品受版权保护。保留一切权利，无论是材料的全部还是部分，特别是翻译、重印、插图使用、朗诵、广播、以缩微胶片或其他方式的复制以及存入数据库的权利。任何形式的使用均须取得版权所有人的许可。

© 2005 Birkhäuser Verlag, P.O. Box 133, CH-4010 Basel, Switzerland

隶属于 Springer Science+Business Media

封面设计：Micha Lotrovsky, 4106 Therwil, Switzerland

印刷用纸为无酸纸，由无氯纸浆制成。TCF ∞

德国印刷

ISBN 3-7643-7153-6

987654321

www.birkhauser.ch

\chapter*{译者的话}
\markboth{译者的话}{译者的话}
\addcontentsline{toc}{chapter}{译者的话}

本书是 Herbert Amann 与 Joachim Escher 所著 \textit{Analysis I}（德文原版
\textit{Analysis I}，Birkhäuser Verlag 1998；英译者 Gary Brookfield，
Birkhäuser Verlag 2005）的中文译本。

\subsection*{关于这个译本}

这是一项\textbf{个人自发的翻译工程}，起因很简单：这本书写得好，而中文读者读它
要跨过一道语言的坎。原书把数系的构造、线性代数的准备与分析的正式展开安排得
极为严整——从 Peano 公理出发依次造出自然数、整数、有理数、实数、复数，再以
度量和范数为基础进入收敛、连续、微分与函数列。这样一条完整的链条，值得有一份
尽可能贴近原文的中文本来对照阅读。

翻译以\textbf{英译本}为底本，同时参照德文原版的书名与体例。全书按章、节组织，
与原文逐节对应，公式、图表、真值表与习题全部保留；章节编号、定理编号
（「命题 1.5」）与公式编号（「(1.4)」）一律沿用原书，以便与英文版、德文版
对照查阅。

\subsection*{关于本译本的生成方式（请先读这一段）}

必须把话说明白：\textbf{本译本由 AI（大语言模型）辅助完成}。原文的识别、初译、
排版代码，以及相当一部分核对工作，都是在人机协作下做的——机器出草稿，人定口径、
抽查扫描件、逐轮返工。它\textbf{不是}由某位译者逐字推敲多年的职业译本，也\textbf{没有经过出版社级别的编校}。

因此请读者务必注意：

\textbf{一、译文可能存在错误。}包括但不限于术语取舍、句子理解、公式与符号的抄录、
编号与交叉引用，以及个别地方与原书不一致的表述。

\textbf{二、不能保证与原著完全一致。}虽已尽量以原书扫描件逐页核对、把发现的讹误
逐条修掉，但「查过」不等于「查全」，遗漏在所难免。

\textbf{三、请勿把本译本当作学术引用、考试或教学的依据。}凡涉及定义、定理与证明的
准确表述，一律以 Amann--Escher 原著（英译本或德文原版）为准。

一句话：可以拿它当\textbf{对照阅读的拐杖}，但不能拿它当\textbf{标准答案}。

\subsection*{译名与记号}

数学术语以国内通行译法为准，重要的名词首次出现时保留英文原词。全书的术语取舍
与记号约定另有对照表，不在此处赘述。原文中以粗体标示的定义性词（如
\textbf{convex}、\textbf{order complete}）在译文中同样以粗体标出。

\subsection*{关于校订}

原书本身有若干印刷瑕疵，扫描件的 OCR 亦不可能全无讹误。凡遇可疑之处，
均以原书扫描件逐页核对为准：属于原书排印问题的，按原样保留而不擅自「改正」；
属于识别错误的，按原书修正。因此本译本中个别地方的写法可能与读者的预期不同，
那通常正是原书的样子。

\clearpage
\subsection*{内容问题反馈}

错误既然难免，就希望它被发现、被记下来。译者建了一份在线反馈表「内容问题反馈表」，
用来收集勘误、术语建议与排版问题——发现哪里不对，随手填一下即可：

\begin{center}
\url{https://github.com/JeaESenna/Amann-Analysis-zh/issues}
\end{center}


反馈时若能附上\textbf{章号、页码与原文对照}，处理起来会快得多。谢谢每一位愿意
回头指正的读者。

\subsection*{免责与致谢}

原著版权归 Herbert Amann、Joachim Escher 与 Birkhäuser Verlag 所有，英译版权
归 Gary Brookfield 所有。本译本仅供个人学习与研究使用，\textbf{不得用于任何商业
用途}。若本书对你有用，请务必去读原著——那是最好的感谢方式。

\vspace{2ex}
\begin{flushright}
Jea神纳\\
2026 年 10 月
\end{flushright}
\chapter*{前言}
\markboth{前言}{前言}
\addcontentsline{toc}{chapter}{前言}

逻辑思维、对复杂关系的分析、对众多问题所共有的深层简单结构的认识——这些是做数学所需的技能，而培养这些技能是数学教育的主要目标。

当然，这些技能不可能“在真空中”习得。唯有坚持不懈地钻研具体问题、孜孜以求深刻的理解，方能走向成功。需要适度的抽象，使人得以专注于本质，而不为表象与枝节所扰。

本书力求清晰与透彻。从一开始，它就要求读者有处理抽象概念的意愿，以及相当程度的主动性。对于这些努力，读者的数学思维能力将得到丰厚的回报，并将具备更深入钻研数学及其应用所必需的基础。

本书是三卷本分析学导论的第一卷。它源自作者们在过去二十六年里于波鸿大学、基尔大学、苏黎世大学、巴塞尔大学和卡塞尔大学讲授的课程。由于我们希望本书也能用于自学和补充阅读，我们收录的材料远多于三个学期所能讲授的内容。这使得我们能对这一学科作广泛的概览，并展示理论诸多优美而重要的应用。我们还表明，数学不仅具有优雅与内在之美，而且为解决具体问题提供了有效的方法。

分析学本身始于第 II 章。在第 I 章中，我们相当透彻地讨论了数系的构造，并介绍了线性代数的基础。这一章尤其适合自学，并提供了从简单假设出发逻辑推导定理的练习。此处，关键在于专注于给定情形中的本质，并避免作出无根据的假设。经验丰富的教师可以轻松地从这一章中选取合适的材料来组成一门课程，也可以在后续各节的学习中按需使用这些基础材料。

在本书中，我们试图为分析学奠定坚实的基础，使读者能够在日后涉足现代数学时在此基础上继续前进。因此，大多数

概念和定义从一开始就以它们的一般形式呈现——即在后续研究和应用中采用的形式。这样，读者只需学习每个概念一次，然后以此为基础，即可直接进入更高等的数学。

我们不打算在此详细描述这三卷的内容，而是请读者参阅各章的导言以及详细的目录。我们还希望提请读者注意每节末尾的众多习题。做这些习题是透彻理解材料的绝对必要条件，同时也是对读者数学进展的有效检验。

在撰写这第一卷的过程中，我们得益于众多同事和学生的建设性批评。我们要特别感谢 Peter Gabriel、Patrick Guidotti、Stephan Maier、Sandro Merino、Frank Weber、Bea Wollenmann、Bruno Scarpellini，以及同样重要的我们的学生，他们以积极的反馈和后来的成就，激励了我们教授分析学的独特方法。

我们从 Peter Gabriel 那里得到了“超出职责范围”的支持。他撰写了附录“数学逻辑导引”，并无私地允许将它收入本书。为此，我们对他深表感谢。

像往常一样，本书成功所需的大量工作是在“幕后”完成的。我们那位“排版完美主义者”的贡献是无法估量的，她在电脑屏幕前度过了无数个小时，并参与了许多关于语法细节的深入讨论。本书的排版与版式完全归功于她，她赢得了我们最诚挚的感谢。

我们还要感谢 Andreas，他为我们提供了最新版本的  \TeX{}\footnote{本书正文由 \LaTeX{} 排版，插图由 CorelDRAW! 与 Maple 绘制。} 并随时准备帮助解决软件和硬件问题。

最后，我们感谢 Thomas Hintermann 鼓励我们将讲义呈献给更广大的读者，也感谢 Thomas Hintermann 与 Birkhäuser Verlag 与我们进行了非常愉快的合作。

苏黎世与卡塞尔，1998 年 6 月

H. Amann 与 J. Escher

\subsection*{第二版前言}

在这一新版中，我们消除了细心的读者向我们指出的错误与不够准确的语言。尤其宝贵的是我们同事 H. Crauel 与 A. Ilchmann 的意见和建议。我们对他们所有人表示由衷的感谢。

苏黎世与汉诺威，2002 年 3 月

H. Amann 与 J. Escher

\subsection*{英译本前言}

我们非常高兴向 Gary Brookfield 致谢，感谢他为将本书译成英文所做的工作。他不仅能够保留德文原文的“精神”，还指出了原版中的不准确之处，并在某些地方提出了更简洁、更清晰的证明，从而帮助改进了数学内容。

苏黎世与汉诺威，2004 年 5 月
